\chapter{Conclusioni}
\label{ch:conclusioni}

\section{Sintesi del lavoro}
\label{sec:sintesi_lavoro}

Il lavoro presentato in questa tesi ha affrontato il problema della generazione automatica di mesh 3D a partire da planimetrie architettoniche CAD, con l'obiettivo di ridurre l'intervento manuale tipicamente richiesto per ottenere una rappresentazione 3D di una scena a partire dal relativo disegno 2D.

È stata realizzata una pipeline in grado di: interpretare le entità geometriche eterogenee di un file DXF, riconducendole a muri, porte e finestre; ricostruire la topologia della planimetria; generare una mesh 3D triangolata, completa di normali e coordinate texture; ed esportare il risultato in formato GLTF, pronto per essere arricchito ed elaborato in Blender ai fini del rendering fotorealistico.

La validazione su diverse planimetrie reali ha mostrato come la pipeline sia in grado di produrre, in tempi dell'ordine di pochi secondi, una ricostruzione geometrica corretta a fronte di un intervento manuale contenuto e circoscritto alla fase di preprocessing, la cui entità dipende dalla conformità del modello di input alle assunzioni progettuali definite in Sezione \ref{sec:assunzioni_progettuali}.

\section{Limiti}
\label{sec:limiti}

L'approccio presenta alcuni limiti, in parte intrinseci alla natura del problema affrontato e in parte legati a scelte progettuali specifiche.

\begin{itemize}
	\item \textbf{Dipendenza dal preprocessing manuale.} La pipeline richiede che il modello di input rispetti dell assunzioni progettuali (Sezione \ref{sec:assunzioni_progettuali}); quando queste non sono soddisfatte, la correzione del modello CAD resta un'operazione manuale.
	
	\item \textbf{Sensibilità ai parametri.} La qualità del risultato dipende dalla corretta taratura di alcuni parametri, come l'unità di misura, la tolleranza $\epsilon$ per il collasso dei vertici (Sezione \ref{sec:collasso_dei_vertici}) e i parametri di DBSCAN per il riconoscimento delle finestre (Sezione \ref{sec:ricostruzione_varchi}); tali parametri non sono al momento determinati automaticamente, ma vanno impostati per ciascun modello.
	
	\item \textbf{Copertura geometrica limitata.} La pipeline è stata progettata e validata su planimetrie composte da muri rettilinei disposti su un unico piano; geometrie più complesse, come muri curvi o edifici multipiano, non sono attualmente supportate.
	
	\item \textbf{Composizione della scena non automatizzata.} Il posizionamento di camera e luci, così come l'eventuale arredamento della scena, sono attualmente operazioni manuali svolte tramite GUI di Blender (Capitolo \ref{ch:rendering}), e non fanno parte della pipeline.
\end{itemize}

\section{Sviluppi futuri}
\label{sec:sviluppi_futuri}

I limiti individuati suggeriscono diverse direzioni di sviluppo futuro. Una prima estensione riguarda l'automazione, anche solo parziale, della fase di preprocessing, ad esempio tramite euristiche in grado di individuare e correggere automaticamente contorni aperti o discontinuità nei muri. 

Sul fronte del rendering, il posizionamento della camera e delle luci potrebbe essere automatizzato tramite euristiche basate sulla geometria della stanza (ad esempio posizionando le luci in funzione della posizione delle finestre), così come l'inserimento di arredi selezionati e disposti automaticamente in base al tipo di stanza.

Infine, un'estensione naturale del lavoro riguarda l'ampliamento della copertura geometrica della pipeline, per supportare planimetrie con muri curvi o edifici composti da più piani, oggi esclusi dalle assunzioni progettuali su cui il sistema si basa.
